\documentclass[12pt]{article}
\usepackage[utf8]{inputenc}
\usepackage[T1]{fontenc}
\usepackage{amsmath,amssymb}
\usepackage{geometry}
\geometry{margin=2.5cm}

\begin{document}

\noindent\underline{\textbf{Corollary}}

\bigskip
\noindent $N\trianglelefteq G \Rightarrow R_N$ is a congruence.

\bigskip
\noindent\underline{Def}: If $G$ is a group and $N\trianglelefteq G \Rightarrow$ the group $G/R_N \overset{\text{not}}{=} G/N$ is called

\noindent the factor (quotient) group of $G$ by the normal subgroup $N$.

\bigskip
\noindent $A\neq\varnothing$, \ $\varphi: A\times A\to A$ a composition law on $A$.

\noindent $R$ = an equivalence relation on $A$ compatible with $\varphi$, \quad $xRx'$ and $yRy' \Rightarrow \varphi(x,y)R\varphi(x',y')$

\noindent = a congruence on $A$.

\[
\frac{A}{R} = \{\hat a \mid a\in A\}, \qquad \varphi^*: \frac{A}{R}\times\frac{A}{R} \to \frac{A}{R}:
\]
\[
\varphi^*(\hat a,\hat b) = \varphi(a,b) \qquad \hat a\hat b=\widehat{ab} \ \text{(multiplicative notation)}
\]

\bigskip
\noindent $\varphi$ associative, commutative $\Rightarrow \varphi^*$ associative, commutative.

\noindent $\varphi$ has an identity element on $A \Rightarrow \hat e$ is the identity element for $\varphi^*$.

\noindent $a\in A$ has an inverse $a'\in A$ with respect to $\varphi \Rightarrow \hat a\in \dfrac{A}{R}$ has an inverse

\noindent $\hat{a'}\in \dfrac{A}{R}$ with respect to $\varphi^*$.

\bigskip
\noindent Hence, if $(A,\varphi)$ is a group $\Rightarrow \Big(\dfrac{A}{R},\varphi^*\Big)$ is a group.

\bigskip
\noindent on $G/N$: \quad $\hat a\hat b = \widehat{ab}$.

\[
aN\cdot bN = abN.
\]

\noindent $\hat e = N$.

\bigskip
\noindent\underline{the factor groups of the group $(\mathbb{Z},+,0)$} \quad $n\mathbb{Z}\leq \mathbb{Z}$, $n\geq 0$.

\[
\mathbb{Z}_n, \ n\geq 1, \qquad n=0 \Rightarrow \frac{\mathbb{Z}}{0} = 0\mathbb{Z} \cong \mathbb{Z}.
\]

\noindent $S_3 \begin{cases} 6 \text{ subgroups} \\ 3 \text{ normal subgroups.} \end{cases}$

\bigskip
\noindent\underline{Remark}: $N\trianglelefteq G$ and ord $G<\infty$. Then ord $\dfrac{G}{N} = [G:N]$.

\[
[\mathbb{Z}:6\mathbb{Z}] = 6 = |\mathbb{Z}_6| = \left|\frac{\mathbb{Z}}{6\mathbb{Z}}\right|
\]

\end{document}
